\section*{Erd\H{o}s Problem 656: a shifted restricted sumset}

\paragraph{Statement.}
If \(A\subseteq\mathbb N\) has positive upper density, must there
exist an infinite \(B\subseteq A\) and an integer \(t\) such that
\[
 \{b+b':b,b'\in B,\ b\ne b'\}+t\subseteq A?
\]

\paragraph{Claim and attribution.}
Yes. We reconstruct the passage from the established dynamical
recurrence theorem to the infinite sumset, retaining its external
inputs explicitly. The source is Kra, Moreira, Richter, and Robertson,
\emph{A proof of Erd\H{o}s's \(B+B+t\) conjecture},
\url{https://arxiv.org/abs/2206.12377}, published in
Communications of the AMS (2024).
The relevant inputs are Proposition 2.3 and Theorem 1.4(i) in
the cited version; the extraction below is its Theorem 2.2.
See also \url{https://www.erdosproblems.com/656}.

\paragraph{The two external inputs.}
Positive upper density implies positive upper Banach density:
choose intervals \([1,N_j]\), with \(N_j\to\infty\), on which the
density converges to a positive limit. These intervals form a
F{\o}lner sequence, meaning that their relative symmetric
difference with any fixed translate tends to zero.

The correspondence principle in Proposition 2.3 provides an
ergodic measure-preserving homeomorphism \(T\) of a compact metric
space \(X\), a point \(a\), a clopen set \(E\) of positive measure,
and a F{\o}lner sequence along which \(a\) is generic, with the
exact identity
\[
 A=\{n\in\mathbb N:T^na\in E\}.
\]
This particular ergodic correspondence principle is an external
input; genericity along the required sequence is not inferred
merely from compactness.

Theorem 1.4(i) then supplies \(x_1,x_2\in X\), \(t\in\mathbb N\),
and strictly increasing integers \(n_j\) such that
\[
 x_1\in E,\qquad T^tx_2\in E,\qquad
 (T^{n_j}a,T^{n_j}x_1)\longrightarrow(x_1,x_2).
                                                               \tag{1}
\]
This recurrence theorem contains the substantial ergodic part
of the known solution. It is used here without claiming a
new proof of that part.

\paragraph{Extracting infinitely many compatible integers.}
Set \(U=E\) and \(V=T^{-t}E\), both open. Then
\(x_1\in U\), \(x_2\in V\).
Pass to a tail of the sequence in (1) so that \(T^{n_j}a\in U\)
for all its terms.
We choose \(b_1<b_2<\cdots\) from this sequence, maintaining
\[
 T^{b_i}x_1\in V\quad\hbox{for every chosen }i,
 \qquad T^{b_i+b_j}a\in V\quad(i<j).
                                                               \tag{2}
\]
For the first choice, (1) gives a sufficiently late \(b_1\)
with \(T^{b_1}x_1\in V\).
Suppose \(b_1,\ldots,b_r\) have been chosen. The open set
\[
 W=U\cap\bigcap_{i=1}^r T^{-b_i}V
\]
contains \(x_1\), by the induction hypothesis.
The convergence in (1) lets us choose \(b_{r+1}>b_r\) with
\[
 T^{b_{r+1}}a\in W,\qquad T^{b_{r+1}}x_1\in V.
\]
This preserves all conditions in (2).
Every selected integer is in \(A\), and for \(i<j\),
\[
 T^{b_i+b_j+t}a\in E,
 \qquad\text{so}\qquad b_i+b_j+t\in A.
\]
Thus \(B=\{b_1,b_2,\ldots\}\) is the required infinite set.
The argument constructs one infinite set, rather than merely
arbitrarily large finite sets with incompatible shifts.

\paragraph{Two elementary checks on the statement.}
If \(A\) has density one, no dynamical theorem is needed and
one can take \(t=0\). Given finitely many chosen \(b_i\in A\),
the finite intersection
\[
 A\cap\bigcap_i(A-b_i)
\]
still has density one, since its complement is contained in a
finite union of sets of density zero. Choose the next, larger
element from this intersection and continue.
In contrast, the shift cannot always be omitted: if \(A\)
is the positive odd integers, the sum of two of its elements
is even, so no infinite \(B\subseteq A\) works with \(t=0\).
Taking an odd shift removes this parity obstruction.

\paragraph{Scope.}
The full answer is proved relative to the two specified established
dynamical inputs. The infinite extraction, the density-one case,
and the need for a shift are proved explicitly.

\paragraph{Completion estimate.}
\textbf{COMPLETION: 100\%}
